\chapter{The Today page}\label{ch:today}

\ui{Today} is the planner's starting point. It answers three questions at a glance: which period needs my
attention, is anything going wrong, and how loaded is the plant in the coming weeks.

\shot{today}{The Today page of the example plant on Saturday 21 November 2026.}{today}

\begin{callouts}
  \item[1] \textbf{Main navigation.} The pages of the daily cycle at the top, plant set-up below, the
           Advanced group folded at the bottom.
  \item[2] \textbf{Factory.} The plant being worked on and its current set-up version; click to open the
           Factory page.
  \item[3] \textbf{Quick search} (\kbd{Ctrl}\,\kbd{K}).
  \item[4] \textbf{Period card.} The period that needs attention now (see below).
  \item[5] \textbf{Needs attention.} Three tiles that turn amber or red when something requires action.
  \item[6] \textbf{Load, next 6 weeks.} Expected load of each stage per week, computed from the open orders.
  \item[7] \textbf{Language and theme.}
  \item[8] \textbf{System status.} Backend reachable, solver queue, database response time, latest schedule.
\end{callouts}

\section{The period card}

The card picks the most relevant period automatically and labels it:

\begin{center}
\begin{tabularx}{\linewidth}{>{\bfseries}p{4.2cm}X}
\toprule
\rowcolor{white}Label & When it is shown \\
\midrule
Period in progress / Current period & Today falls inside the period. \\
Next period & No period covers today, and this is the next one that starts. \\
Past period not closed yet & The period has ended but was never closed, so its undelivered orders are stuck in it. Close it (Section~\ref{sec:close}). \\
\bottomrule
\end{tabularx}
\end{center}

Under the period name you see its dates and four tiles:

\begin{itemize}
  \item \ui{Orders in period}: orders due inside the period, plus (``+N after'') the orders in the look-ahead
        window that could be pulled ahead;
  \item \ui{Input errors}: problems that block scheduling (and the number of warnings, which do not block);
  \item \ui{Options}: how many feasible schedules exist, and how many are still \ui{solving};
  \item \ui{Committed}: the name of the committed option and when it was committed, or a dash.
\end{itemize}

The button on the right always proposes the next step: \ui{Check the inputs}, \ui{Fix the inputs},
\ui{Generate options}, \ui{Watch the solver}, \ui{Compare \& commit}, \ui{Today's work orders} or
\ui{Open the period}. If the period has what-if variants or reschedules, they are listed as \ui{Other
versions} below the tiles.

When no open period exists the card says \ui{No open period} and offers to create one.

\section{Needs attention}

\begin{description}[style=nextline, leftmargin=1.2em, font=\bfseries\color{brand}]
  \item[Overdue open orders] Orders whose due date has passed but which are still open in the order book.
        Usually they belong to a period that was not closed, or nobody scheduled them. Schedule them as
        early as possible and tell the customer.
  \item[Orders due $\leq$ 7 days, no period] Orders due within a week that no open period contains. Create a
        period or update an existing one from the order book.
  \item[Solver at work] Runs that are queued or solving. The tile links to the Runs page filtered on running
        runs. When a period has input errors, the tile \ui{Periods with input errors} is shown instead.
\end{description}

Each tile is a link to the page where the problem is fixed.

\section{Load, next 6 weeks}

The table shows, for each stage and each of the next six weeks, the machine time the open orders need as a
percentage of the machine time the calendar provides (shifts minus breaks, rest days, holidays and weekly
maintenance; furnaces that run 24/7 count the whole day). The estimate uses the fastest machine and excludes
changeovers, so it is a lower bound. Green cells are comfortable, blue cells are busy, and amber cells above
85~percent are tight. In the example, heat treatment at 64--78~percent is clearly the bottleneck stage.

Hover a cell to see the absolute figures (``need / available machine-hours''). The link \ui{Order book}
opens the same table over eight weeks, together with the orders behind it.

\section{Other open periods and latest runs}

At the bottom, \ui{Other open periods} lists the remaining open periods, and \ui{Latest solver runs} shows
the most recent algorithm runs with their status (\ui{Succeeded}, \ui{Solving}, \ui{Queued}, \ui{Failed},
\ui{Cancelled}). \ui{View all} opens the Runs page (Section~\ref{sec:runs}).
